\section{Theoretical Analysis}
\label{sec:submodular}

This section establishes the two facts the decomposition rests on. Once the
firing weapons are fixed, assigning targets among them greedily by marginal
value attains at least half the stage optimum with no learned parameters
(Proposition~\ref{prop:auction-half}), and the guarantee holds for any firing
set, including one produced by all weapons deciding at the same instant. The
other half of the decision admits no such treatment: choosing which weapons
fire by what is best in the current stage can be arbitrarily far from optimal
(Proposition~\ref{prop:myopic-gap}). That gap is the entire job of the learned
policy.

Throughout we reuse the stage objective $f_t$ of Eq.~\eqref{eq:round-obj},
the firing set $F$, and the edge set $E_t$ introduced in
Section~\ref{sec:auction}, where $f_t$ was shown to be monotone
submodular~\cite{wang2013weapon} under a partition matroid induced by
Eq.~\eqref{eq:weapon_target_constraint}. Write
$\text{OPT}_t(F) = \max_{S_t \text{ feasible}} f_t(S_t)$ and
$\text{OPT}_t = \text{OPT}_t(\{1,\ldots,M\})$.

\subsection{Target Assignment Is Solved, Not Approximated}

%% ============================================================
%% COMMENTED OUT 2026-09-14: Propositions par-linear and seq-half.
%%
%% These two set up the standard case that simultaneous multi-agent decoding
%% needs a coordination mechanism, so that the assignment guarantee could then
%% be shown to remove the need for one. That argument turned out to be longer
%% than it needs to be: prop:auction-half holds for ANY firing set, which
%% already says the guarantee does not care whether the set was produced all
%% at once, and the comparison against sequential decoding is carried
%% empirically by Table~\ref{tab:main_results}.
%%
%% Restore if a reviewer asks why simultaneous decoding is safe in the
%% worst case rather than only on the benchmark. related_works.tex also
%% cited this pair; that sentence was rewritten when these were cut.
%% ============================================================
%% \begin{proposition}[Parallel decoding without an assignment mechanism: tight $\Theta(1/M)$]
%% \label{prop:par-linear}
%% Consider the idealized zero-communication parallel rule in which each weapon
%% $m$ simultaneously selects its own feasible pair of largest singleton value
%% $f_t(\{(m,n)\})$, with no visibility into the other weapons' same-stage
%% choices. Then $f_t(\mathrm{Par}_t) \geq \tfrac{1}{M}\text{OPT}_t$, and this
%% is tight.
%% \end{proposition}
%% \begin{proof}
%% \emph{Lower bound.} Let $e^\ast = \arg\max_{e\in E_t} f_t(\{e\})$. By
%% monotonicity and submodularity $f_t$ is subadditive, so for the optimal
%% feasible $S^\ast_t$ (with $|S^\ast_t|\leq M$),
%% $\text{OPT}_t \leq \sum_{e\in S^\ast_t} f_t(\{e\}) \leq M f_t(\{e^\ast\})$.
%% The weapon owning $e^\ast$ selects it, since $e^\ast$ is globally best and
%% hence best within its own part, so $f_t(\mathrm{Par}_t) \geq f_t(\{e^\ast\})
%% \geq \text{OPT}_t/M$.
%%
%% \emph{Tightness.} Take $M$ weapons and $M{+}1$ targets: a shared target $n_0$
%% of value $1$ and one private target per weapon of value $1-\varepsilon$, all
%% with $P\equiv 1$ and all pairs feasible. Every weapon's best singleton is
%% $n_0$, so $\mathrm{Par}_t$ sends all $M$ weapons to $n_0$, giving
%% $f_t(\mathrm{Par}_t)=1$, while assigning one weapon to $n_0$ and the rest to
%% their private targets gives $\text{OPT}_t = 1+(M{-}1)(1{-}\varepsilon)$. The
%% ratio tends to $1/M$ as $\varepsilon\to 0$.
%% \end{proof}
%%
%% \begin{proposition}[Sequential decoding carries a constant-factor guarantee]
%% \label{prop:seq-half}
%% An idealized sequential rule, in which each weapon in turn selects the
%% feasible pair of largest marginal gain given all pairs already selected this
%% stage, is the greedy algorithm for monotone submodular maximization under a
%% matroid constraint, and therefore guarantees
%% $f_t(\mathrm{Seq}_t) \geq \tfrac12\text{OPT}_t$~\cite{fisher1978analysis}.
%% \end{proposition}
%%
%% Propositions~\ref{prop:par-linear} and~\ref{prop:seq-half} are the standard
%% argument for why simultaneous multi-agent decoding needs a coordination
%% mechanism, and they are why the multi-agent reinforcement learning
%% literature has largely moved to sequential, auto-regressive action
%% generation: the multi-agent advantage decomposition
%% theorem~\cite{kuba2022trust} transforms the joint policy search into exactly
%% the sequential decision process of Proposition~\ref{prop:seq-half}, and it is
%% that sequencing which buys the monotonic improvement guarantees of
%% HATRPO/HAPPO~\cite{kuba2022trust} and MAT~\cite{wen2022mat}.
%%

Our pipeline does not coordinate the weapons' target choices better. It
removes the choice from them.

\begin{proposition}[Many-to-one marginal assignment is matroid greedy on the stage objective]
\label{prop:auction-half}
Fix stage $t$ and a firing set $F$, and let $\text{OPT}_t(F)$ be the maximum of $f_t$ over feasible allocations of the weapons in $F$. The many-to-one greedy assignment returns an allocation $S^{\text{gre}}_t$ with
\begin{equation}
f_t(S^{\text{auc}}_t) \;\geq\; \tfrac12\,\text{OPT}_t(F).
\end{equation}
\end{proposition}
\begin{proof}
First, the assignment score equals the exact marginal gain of $f_t$. Let $S$ be the allocation so far and $S_{n'}$ the weapons already assigned to target $n'$. Adding $(m',n')$ changes only the $n'$ summand of Eq.~\eqref{eq:round-obj}:
\begin{align*}
f_t(S \cup \{(m',n')\}) - f_t(S)
&= V_{n'}\rho_{n'}^{(t)}\Big[\textstyle\prod_{m \in S_{n'}}(1-P_{m,n'}) \\
&\qquad\qquad - \textstyle\prod_{m \in S_{n'} \cup \{m'\}}(1-P_{m,n'})\Big] \\
&= V_{n'}\rho_{n'}^{(t)} \Big(\textstyle\prod_{m \in S_{n'}}(1-P_{m,n'})\Big)\big[1-(1-P_{m',n'})\big] \\
&= V_{n'}\rho_{n'}^{(t)}\,\sigma_{n'}\,P_{m',n'} \;=\; b_{m',n'},
\end{align*}
which is the quantity of Eq.~\eqref{eq:marginal-bid}. Second, the rule is greedy in that quantity: each iteration takes the single largest score over all unassigned weapons and legal targets, and never revokes an assignment. It is therefore the greedy algorithm for maximizing the monotone submodular $f_t$ under the partition matroid constraint, which attains at least $\tfrac12$ of the constrained optimum~\cite{fisher1978analysis}.
\end{proof}

The factor $\tfrac12$ is a worst-case floor that holds regardless of the scale of the problem: it does not degrade as the number of weapons $M$ grows. It is also close to the best attainable, since no polynomial-time method can guarantee more than $1-1/e$ for this problem class~\cite{feige1998threshold}, and the alternatives analyzed in Section~\ref{sec:submodular} carry either no guarantee or one that decays as $1/M$. The floor is also loose in practice. On the $(5,5,5)$ and $(5,7,5)$ benchmark configurations, where $\text{OPT}_t(F)$ can be computed exactly by enumerating every feasible allocation of the firing weapons, the rule attains a mean 99.3\% of the stage optimum over 77 policy-generated stages (ten instances per configuration, distribution of Section~\ref{sec:data}) and never falls below 87.1\%. The $\tfrac12$ bound is thus a guarantee against adversarial instances, not a description of typical behavior.
%% Measured 2026-08-24 by code_DWTA/eval_auction_stage_ratio.py (result/auction_stage_ratio.txt):
%% (5,5,5) mean 0.9939 min 0.8711 over 38 stages; (5,7,5) 0.9918/0.8782 over 39; combined 77 stages
%% mean 0.9928 min 0.8711. A third probe config (5,15,5) gave 0.9973/0.9496 (43 stages) -- not reported
%% in prose per author decision (non-benchmark config). Ratio = f_t(auction)/OPT_t(F) with F fixed to
%% the trained SCoPE-Comm policy's firing set, OPT by exhaustive enumeration. Stages with empty F or
%% no legal target excluded.

What matters for the rest of the argument is the quantifier in
Proposition~\ref{prop:auction-half}, not the constant.

\begin{remark}
\label{rem:auction-kills-pileon}
Proposition~\ref{prop:auction-half} holds for \emph{any} firing set, in
particular one produced by $M$ weapons deciding simultaneously and
independently. Weapons deciding at the same instant would normally contend
for the same target, and a communication layer or conflict handler exists to
settle that contention. Here there is nothing to settle: no weapon chooses a
target, so the choices that would collide are never made. The guarantee is
obtained at one network call per stage rather than $M$.
\end{remark}

\subsection{Participation Cannot Be Decided Per Stage}

Proposition~\ref{prop:auction-half} is conditional on $F$. It says nothing
about how $F$ should be chosen, and the natural per-stage answer is to fire
every weapon whose best available marginal value is positive. That answer is
unboundedly bad.

\begin{proposition}[Myopic participation is arbitrarily suboptimal]
\label{prop:myopic-gap}
Let $\textsc{Myopic}$ be the rule that, at every stage, fires every weapon
with at least one legal target of positive marginal value, and assigns
targets by greedy marginal assignment. For every $\kappa > 0$ there is a DWTA instance on
which the value destroyed by $\textsc{Myopic}$ is less than $\kappa$ times
the value destroyed by an optimal policy. No constant-factor guarantee for
$\textsc{Myopic}$ exists.
\end{proposition}
\begin{proof}
Take $T=2$, $M=1$, $N=2$. The single weapon has ammunition $A_1 = 1$, with
$P_{1,1}=P_{1,2}=p \in (0,1]$. Target 1
has value $V_1 = 1$ and time window $TW_1 = \{1\}$; target 2 has value
$V_2 = V > 0$ and time window $TW_2 = \{2\}$.

At stage 1 the only legal engagement is target 1, whose marginal value
$V_1 p = p$ is positive, so $\textsc{Myopic}$ fires and exhausts the
weapon's ammunition. At stage 2 target 2 is legal but the weapon has no
ammunition, so nothing is engaged. $\textsc{Myopic}$ destroys $p$.

The policy that holds at stage 1 and fires at target 2 at stage 2 is
feasible and destroys $pV$. Hence the optimum destroys at least $pV$, and
the ratio is at most $p/(pV) = 1/V$. Choosing $V > 1/\kappa$ gives the
claim.
\end{proof}

\begin{remark}
\label{rem:myopic-is-the-job}
The instance in Proposition~\ref{prop:myopic-gap} isolates the one quantity
no within-stage mechanism can access: the value of \emph{not} acting, which
is zero within the stage and realized only later. Both the assignment
guarantee (Proposition~\ref{prop:auction-half}) is per-stage, and so says
nothing about it. This is the precise sense in which the
decomposition of Section~\ref{sec:method} assigns learning only the part
that cannot be obtained otherwise. It is also directly measurable: the
Greedy baseline of Table~\ref{tab:main_results} is $\textsc{Myopic}$: it uses
the same assignment rule as our pipeline and fires every weapon it can, so
the margin between it and the learned pipeline is an empirical estimate of
how much this gap costs on realistic instances rather than adversarial ones.
\end{remark}

\begin{remark}
\label{rem:gap-interpretation}
Proposition~\ref{prop:auction-half} concerns an idealized rule, not a
trained network, and bounds a single stage rather than a full $T$-stage
trajectory. Value missed at one stage is not necessarily
lost: it remains on the battlefield, is re-observed through the updated
$\rho_n^{(t+1)}$, and can be recovered later, which is one reason measured
gaps are far smaller than the worst cases above. Proposition~\ref{prop:myopic-gap}
is the exception, and deliberately so: its construction uses a closing time
window precisely so that the missed value cannot be recovered, which is what
makes the participation decision genuinely multi-stage rather than merely
harder-to-see.
\end{remark}
